% Table 5: explanatory gain along the ladder
% GENERATED by src/build_tables.py from data/numbers.json.
% Do not edit by hand; rerun the generator.
\begin{table}[pos=htbp]
\tableserif
\centering
\begin{threeparttable}
\hllabelcolour{HlM}
\caption{Scale-robust \hlmulti{explanatory gain from the structured to the} narrative-augmented \hlmulti{specification}.}
\label{tab:ladder}
\small
\renewcommand{\arraystretch}{1.2}
\begin{tabularx}{\textwidth}{>{\raggedright\arraybackslash}X ccc}
\toprule
\headerrow
\hcell{Scale-robust metric} & \hcell{M0} & \hcell{M0c} & \hcell{M1} \\
\midrule
Adjusted $\rho^2$ (McFadden) & \textbf{0.0150} & 0.0109 & 0.0002 \\
\altrow
Cross-validated log loss (SD) & \textbf{0.829 (0.053)} & 0.836 (0.055) & 0.849 (0.056) \\
AIC & \textbf{849.1} & 852.7 & 861.9 \\
\altrow
BIC & \textbf{976.8} & 988.9 & 1019.4 \\
Events per parameter & 2.57 & 2.41 & 2.08 \\
\altrow
Log-likelihood $\log L$ & $-$394.55 & $-$394.33 & $-$393.93 \\
Covariates / parameters & 28 / 30 & 30 / 32 & 35 / 37 \\
\addlinespace
\rowcolor{PanelBg}
\multicolumn{4}{l}{\textbf{Likelihood-ratio tests between adjacent rungs}} \\
\midrule
\headerrow
\hcell{Comparison} & \hcell{$\chi^2$} & \hcell{d.f.} & \hcell{$p$} \\
\midrule
M0 $\rightarrow$ M0c, coded contributing factors & 0.44 & 2 & 0.803 \\
\altrow
M0c $\rightarrow$ M1, narrative-derived cues & 0.80 & 5 & 0.977 \\
\bottomrule
\end{tabularx}
\begin{tablenotes}[flushleft]\footnotesize
\item[] \textit{Note:} The ladder is the sequence of nested specifications, and each
specification on it is a rung. M0 holds the screened structured covariates, M0c adds the
officer-coded contributing factors, and M1 adds the narrative-derived pre-crash cues.
Adjusted $\rho^2$ penalizes the McFadden $\rho^2$ for the additional parameters of each
rung. Cross-validated log loss is the mean over 10 repeats of
5-fold stratified cross-validation, with the standard deviation across
folds in parentheses. AIC and BIC add to $-2\log L$ a penalty of 2 and of $\ln n$ per
parameter, with $n$ the number of riders. Bold marks the best value in each of the first four rows, the highest adjusted
$\rho^2$ and the lowest log loss, AIC and BIC. Fold-level values and the paired differences
between rungs are in Supplementary Table~\suppTabCvfolds{}. A large $p$ in a likelihood-ratio test
means that no gain was detected at this sample size, not that no gain exists.
\end{tablenotes}
\end{threeparttable}
\end{table}
